
Reproducibility badges have become standard at major machine learning venues, requiring code and data artifact deposition. Despite five years of badge programs, their effectiveness remains empirically unverified. This study provides the first continuous quality-outcome measurement linking documentation artifact quality to reproducibility in ML benchmarks. We analyzed 108 classification benchmarks (2019-2024) with $\geq 5$ independent reproduction attempts each, using performance variance (coefficient of variation) across independent attempts as a scalable reproducibility proxy. Mean artifact quality scored 2.43/10 (threshold: 7.0), with critical gaps in evaluation protocols (1.19/10) and hyperparameters (1.16/10). Automated scoring consistency was perfect ($\kappa$=1.0), confirming measurement reliability. Despite artifact presence at scale, performance variance showed no statistically significant reduction: high-artifact benchmarks ($\geq 2$ artifacts) exhibited mean CV=0.035 versus low-artifact benchmarks mean CV=0.069 (Mann-Whitney p=0.418, Cohen's d=0.464). While the directional trend is positive, the effect size falls below the medium threshold (0.5) and the sample (n=22 benchmarks in variance analysis) was underpowered relative to target (n=100). Our findings suggest reproducibility badges increase artifact presence but not quality, producing no detectable reproducibility benefit. We conclude that badge programs require quality enforcement mechanisms---rubric-based scoring, post-publication audits, or automated completeness checks---to fulfill their intended purpose.
